\scanpage{41}
\[
H^*(X)_T=H^*(X)\otimes_K\widehat K_T.
\]
une décomposition,
\[
\boxed{\varphi_T:\coprod_{p,q}H^q(X,\Omega_X^p)_T\longrightarrow H^*(X)_T}
\]

Cela permet de définir
\[
H_T^{p,q}(X)=\varphi_T\big(H^q(X,\Omega_X^p)_T\big)\cap H^*(X)
 \subset F^p\big(H^q(X)\big)
\]
qui est un \uline{sous}-\inserted{$K$-}\uline{espace vectoriel de} $H^*(X)$,
défini canoniquement par la donnée de $T$. Nous
\unsure{poserons|posons|pouvons} de $=$
\[
H_T^{pq}(X,\mathbb R)=\varphi_T\big(H^q(X,\Omega_X^p)_T\big)
\cap H^*(X_{\widehat K_T},\mathbb R)\cap H^*(X)
\]
\deleted{c'est un \unsure{sous} \unsure{$K$} \unsure{$\mathbb R\cap K_T$}-espace vectoriel de}
c'est encore un sous-\inserted{$K$-}espace vectoriel de $H^*(X)$
si $\widehat K_T\simeq\mathbb R$ ; si\deleted{non} \inserted{$\widehat K_T\simeq\mathbb C$} c'est un sous-\deleted{\unsure{$R_{\mathbb R}$|$\mathbb R_T$}}$R_T$-espace vectoriel, où
\[
R_T=\mathbb R\cap K\qquad\text{dans }\widehat K_T\simeq\mathbb C
\]
[sous-\deleted{\unsure{espace}}\inserted{corps} \uline{réel} pour $T$] ; à fortiori, \deleted{\unsure{c'est}} dans tous les
cas, on obtient un $K_r^{\unsure{*}}$-espace \unsure{vectoriel|vectorel}, où
\[
K_r\simeq\bigcap_T R_T\simeq\text{plus \unsure{grand|gd} sous-corps tot. réel de }K
\]
(NB $K_r$ est une ext. algébrique du corps premier $\mathbb Q$).
% archive sheet number: 40
